\section{Related Work}
\label{sec:related_work}

\paragraph{Reinforcement Learning with Verifiable Rewards.}
%Reinforcement learning from human feedback aligns language models with preferences through learned reward models~\citep{christiano2017deep,ouyang2022training,bai2022training}, with DPO offering a verifier-free preference alternative~\citep{rafailov2023direct}. 
Reinforcement learning with verifiable rewards (RLVR) instead replaces preference models with programmatic correctness checks on tasks whose answers are objectively verifiable~\citep{shao2024deepseekmath,deepseekr1,lambert2024tulu3,zeng2503simplerlzoo}. 
Recent work refines GRPO-style optimization through dynamic sampling and clipping~\citep{yu2025dapo}, extended training schedules~\citep{liu2025prorl}, finer credit assignment~\citep{kazemnejad2025vineppo}, asymmetric handling of negative samples~\citep{zhu2025negsample}, entropy regularization~\citep{cui2025entropy}, and simplified baselines~\citep{ahmadian2024back,hu2025reinforce++}. Outcome-only signals nevertheless remain structurally blind: traces with very different dependency support can be reward-equivalent~\citep{chen2026does}, and reward hacking persists~\citep{gao2023scaling,skalse2022defining}. \textit{\method injects deterministic topology and continuity diagnostics from latent reasoning DAGs, targeting structural credit assignment directly rather than through learned reward surrogates.}

\paragraph{Process Reward Models for Reasoning.}
% Process reward models (PRMs) score intermediate reasoning steps rather than only final answers, with early variants relying on dense human step annotations~\citep{lightman2023lets}.
% Later work reduces annotation cost with Monte Carlo or tree-search rollouts~\citep{wang2024math,luo2024omegaprm,zhang2025qwen25mathprm}, while generative PRMs model verifier reasoning explicitly~\citep{thinkprm2025,genprm2026,she2025rprm}. Implicit PRM~\citep{yuan2024implicitprm} derives step-level signals from outcome feedback, yet PRMBench~\citep{song2025prmbench} shows persistent weaknesses on fine-grained error categories. TopoPRM is complementary to this line: it does not train a neural verifier, and it does not claim proof-level verification; instead, it introduces deterministic global-structure supervision over dependency topology and local continuity.

Process reward models (PRMs) score intermediate steps to address the sparsity of outcome rewards, originally relying on dense human annotation~\citep{uesato2022solving,lightman2023lets}. Subsequent work reduces labeling cost via Monte Carlo or tree-search rollouts~\citep{wang2024math,luo2024omegaprm,zhang2024rest,lu2024autopsv,zhang2025qwen25mathprm}, generative PRMs make verifier reasoning explicit~\citep{thinkprm2025,genprm2026,she2025rprm}, and process advantage verifiers tie step values to rollout-induced returns~\citep{setlur2025rewarding,lai2024step}. Implicit PRMs derive step-level signals from outcome feedback alone~\citep{yuan2024implicitprm}, yet recent benchmarks expose persistent weaknesses on fine-grained error categories and inter-step consistency~\citep{song2025prmbench,zheng2025processbench}. A shared limitation is that scoring remains step-local, treating each step as conditioned only on its immediate predecessor. \textit{\method trains no neural verifier and claims no proof-level guarantee, instead introducing deterministic global-structure supervision over dependency topology together with local continuity.}


\paragraph{On-Policy Distillation and Reasoning Compression.}
% Knowledge distillation~\citep{hinton2015distilling} compresses large models into smaller students, and on-policy distillation (OPD) closes the train-test gap by letting the student generate its own trajectories under teacher supervision~\citep{minillm2023,agarwal2024gkd,song2026opdsurvey}.
% Recent OPD variants for reasoning include OPSDC~\citep{opsdc2026}, ExOPD~\citep{yang2026exopd}, BOND~\citep{sessa2024bond}, and KDRL~\citep{xu2025kdrl}; black-box alternatives such as GAD~\citep{ye2025gad} further broaden applicability. However, most methods still optimize token-level matching objectives without explicit structural control over dependency preservation. In parallel, interpretability studies show graph-like latent organization in long-CoT reasoning~\citep{besta2024graph,zhong2026chains2dags}. TopoPRM bridges these lines by using one deterministic DAG diagnostic interface for both reward optimization and on-policy distillation, so compression is conditioned on structural adequacy rather than only token imitation.

Knowledge distillation transfers capability from larger teachers~\citep{hinton2015distilling}, and on-policy variants close the train-test gap by supervising student-generated trajectories~\citep{minillm2023,agarwal2024gkd,ko2024distillm,song2026opdsurvey}. Reasoning-oriented OPD includes OPSDC~\citep{opsdc2026}, ExOPD~\citep{yang2026exopd}, BOND~\citep{sessa2024bond}, and KDRL~\citep{xu2025kdrl}, with black-box GAD broadening applicability beyond logits access~\citep{ye2025gad}; in parallel, long-CoT compression has been pursued through token pruning, length penalties, and explicit budget control~\citep{xia2025tokenskip,luo2025o1,arora2026training,aggarwal2025l1,muennighoff2025s1}. Most of these objectives still optimize token-level matching without explicit structural constraints, even as interpretability studies reveal graph-like latent organization in long-CoT reasoning~\citep{yao2023tree,besta2024graph,zhong2026chains2dags}. \textit{\method drives both topology-aware reward optimization and on-policy self-distillation, so that compression is conditioned on structural adequacy rather than token imitation alone.}
